\subsection{Localization of homological dimension}

PROPOSITION 2.--- Let A be a ring, M and N be A-modules, i an integer, and S a multiplicative subset of \(A\). We have a canonical isomorphism of \(\mathrm{S}^{-1}\mathrm{A}\) -modules

\[
\mathrm{S} ^ {- 1} \operatorname{Tor} _ {i} ^ {A} (\mathrm{M}, \mathrm{N}) \longrightarrow \operatorname{Tor} _ {i} ^ {\mathrm{S} ^ {- 1} A} (\mathrm{S} ^ {- 1} \mathrm{M}, \mathrm{S} ^ {- 1} \mathrm{N}).
\]

If the ring A is Noetherian and the A-module M is finitely generated, one has a canonical isomorphism of S \(^{-1}\) A-modules

\[
\mathrm{S} ^ {- 1} \operatorname{Ext} _ {A} ^ {i} (\mathrm{M}, \mathrm{N}) \longrightarrow \operatorname{Ext} _ {\mathrm{S} ^ {- 1} \mathrm{A}} ^ {i} (\mathrm{S} ^ {- 1} \mathrm{M}, \mathrm{S} ^ {- 1} \mathrm{N}).
\]

As the A-module \(S^{-1}A\) is flat, this follows from A, X, p.~110, prop. 9 and p.~111, prop. 10.

COROLLARY.--- Let A be a ring, M and N be A-modules, and i an integer.

\begin{enumerate}
\def\labelenumi{\alph{enumi})}
\item
  The support of \(\mathrm{Tor}_{i}^{A}(\mathrm{M},\mathrm{N})\) is contained in \(\mathrm{Supp}(\mathrm{M})\cap \mathrm{Supp}(\mathrm{N})\) and the same is true of the support of \(\mathrm{Ext}_{\mathrm{A}}^{i}(\mathrm{M},\mathrm{N})\) if A is Noetherian and M is finitely generated.
\item
  Suppose A is Noetherian and the modules M and N are finitely generated; if the A-module M⊗A N has finite length, then so do Tor \(_{i}^{A}\) (M, N) and Ext \(_{A}^{i}\) (M, N).
\end{enumerate}

If p is a prime ideal of A not belonging to Supp(M) ∩ Supp(N), one of the modules \(M_{p}\) or \(N_{p}\) is zero, which implies a) in view of prop. 2.

For a finitely generated module over a Noetherian ring to be of finite length, it is necessary and sufficient that its support consist of maximal ideals (IV, § 2, no. 5, prop. 7). Under hypothesis b), the A-modules Tor \(_{i}^{A}\) (M, N) and Ext \(_{A}^{i}\) (M,N) are finitely generated (A, X, p.~108, cor.); assertion b) therefore follows from a).

PROPOSITION 3.--- Let A be a Noetherian ring, M a finitely generated A-module, and N an A-module.

\begin{enumerate}
\def\labelenumi{\alph{enumi})}
\tightlist
\item
  We have
\end{enumerate}

\[
\mathrm{dp} _ {\mathrm{A}} (\mathrm{M}) = \sup _ {\mathfrak {p}} \mathrm{dp} _ {\mathrm{A} _ {\mathfrak {p}}} \left(\mathrm{M} _ {\mathfrak {p}}\right), \quad \mathrm{di} _ {\mathrm{A}} (\mathrm{N}) = \sup _ {\mathfrak {p}} \mathrm{di} _ {\mathrm{A} _ {\mathfrak {p}}} \left(\mathrm{N} _ {\mathfrak {p}}\right), \quad \mathrm{dh} (\mathrm{A}) = \sup _ {\mathfrak {p}} \mathrm{dh} \left(\mathrm{A} _ {\mathfrak {p}}\right),
\]

where p ranges over the set of prime (resp. maximal) ideals of A.

\begin{enumerate}
\def\labelenumi{\alph{enumi})}
\setcounter{enumi}{1}
\tightlist
\item
  The map \(\mathfrak{p} \mapsto \mathrm{dp}_{A_{\mathfrak{p}}}(\mathrm{M}_{\mathfrak{p}})\) of \(\operatorname{Spec}(A)\) in \(\overline{Z}\) is upper semi-continuous.
\end{enumerate}

Let us prove a). Let \(n\) be an integer \(\geqslant 0\). Suppose we have \(\mathrm{dp}_{A}(\mathbf{M}) < n\). For every prime ideal \(\mathfrak{p}\) of A and every \(\mathrm{A}_{\mathfrak{p}}\) -module Q, the \(\mathrm{A}_{\mathfrak{p}}\) -module \(\mathrm{Ext}_{\mathrm{A}_{\mathfrak{p}}}^{n}(\mathrm{M}_{\mathfrak{p}},\mathrm{Q})\) is isomorphic to \((\mathrm{Ext}_{\mathrm{A}}^{n}(\mathrm{M},\mathrm{Q}))_{\mathfrak{p}}\) (prop. 2), hence is zero; from this one deduces the inequality \(\mathrm{dp}_{A_{\mathfrak{p}}}(\mathrm{M}_{\mathfrak{p}}) \leqslant \mathrm{dp}_{\mathrm{A}}(\mathrm{M})\) (A, X, p.~134, prop. 1). Suppose conversely that one has \(\mathrm{dp}_{A_{\mathfrak{m}}}(\mathrm{M}_{\mathfrak{m}}) < n\) for every maximal ideal \(\mathfrak{m}\) of A, and let R be an A-module. One has \((\operatorname{Ext}_{\mathrm{A}}^{n}(\mathrm{M}, \mathrm{R}))_{\mathfrak{m}} = 0\) for every \(\mathfrak{m}\) (prop. 2), hence \(\operatorname{Ext}_{\mathrm{A}}^{n}(\mathrm{M}, \mathrm{R}) = 0\) (II, § 3, n° 3, cor. 2 of th. 1), which implies \(\mathrm{dp}_{\mathrm{A}}(\mathrm{M}) < n\) (A, X, p.~134, prop. 1). The first equality of a) follows from this. The second is proved in the same way, using the characterization of injective dimension given in prop. 1 (iii). Since dh(A) is the upper bound of the set of injective dimensions of A-modules (n° 1), the third equality follows from this.

Let us prove b). Let \(\mathfrak{p}\) be a prime ideal of A and \(n = \mathrm{dp}_{\mathrm{A}_{\mathfrak{p}}}(M_{\mathfrak{p}})\) . Let us show that there exists a neighborhood U of \(\mathfrak{p}\) in \(\operatorname{Spec}(A)\) such that one has \(\mathrm{dp}_{\mathrm{A}_{\mathfrak{q}}}(M_{\mathfrak{q}}) \leqslant n\) for every \(\mathfrak{q} \in U\) . This is clear if \(n = +\infty\) ; if \(n = -\infty\) , this follows from the fact that the support of M is closed. Now suppose \(n\) is finite and choose an exact sequence of A-modules

\[
\mathrm{P} _ {n - 1} \stackrel {d _ {n - 1}} {\longrightarrow} \mathrm{P} _ {n - 2} \longrightarrow \dots \longrightarrow \mathrm{P} _ {0} \stackrel {d _ {0}} {\longrightarrow} \mathrm{M} \rightarrow 0,
\]

where the \(P_{i}\) are free of finite type (A, X, p.~53, prop. 6). Set \(P = Ker d_{n-1}\) ; it is a finitely presented module. The \(A_{p}\) -module \(P_{p}\) is projective (A, X, p. 134, prop. 1), hence free (II, § 3, no. 2, cor. 2 of prop. 5). By II, § 5, no. 1, cor. of prop. 2, there exists an element f of \(A - p\) such that the \(A_{f}\) -module \(P_{f}\) is free; the \(A_{q}\) -module \(P_{q}\) is then free for every prime ideal q of the open set U of Spec(A) consisting of the prime ideals not containing f. This proves b).

COROLLARY 1.--- For every multiplicative subset S of A, we have

\[
\mathrm{dp} _ {\mathrm{S} ^ {- 1} \mathrm{A}} (\mathrm{S} ^ {- 1} \mathrm{M}) \leqslant \mathrm{dp} _ {\mathrm{A}} (\mathrm{M}), \quad \mathrm{di} _ {\mathrm{S} ^ {- 1} \mathrm{A}} (\mathrm{S} ^ {- 1} \mathrm{N}) \leqslant \mathrm{di} _ {\mathrm{A}} (\mathrm{N}), \quad \mathrm{dh} (\mathrm{S} ^ {- 1} \mathrm{A}) \leqslant \mathrm{dh} (\mathrm{A}).
\]

COROLLARY 2.--- If one has dp \(_{A_m}\) (M \(_m\) ) \textless{} +∞ for every maximal ideal m of Supp(M), we have dp \(_A\) (M) \textless{} +∞.

Indeed, the subspace X of Supp(M) consisting of the maximal ideals is quasi-compact (II, § 4, no. 2, props. 8 and 9 and no. 3, cor. 7 of prop. 11); the map \(\mathfrak{m} \mapsto \mathrm{dp}_{\mathrm{A}_{\mathfrak{m}}}(\mathrm{M}_{\mathfrak{m}})\) of X into \(\overline{\mathbf{R}}\) is upper semicontinuous (prop. 3), hence bounded (TG, IV, p.~30, cor. of th. 3).

Remark. --- Let A be a regular Noetherian ring of infinite dimension (VIII, § 5, exerc. 6 c)). We shall see below (n° 7, th. 2 and § 4, n° 1, prop. 1) that one has \(\mathrm{di}_{\mathrm{A_m}}(\mathrm{A_m}) = \mathrm{dh}(\mathrm{A_m}) < +\infty\) for every maximal ideal m of A; Proposition 3 therefore implies \(\mathrm{di}_{\mathrm{A}}(\mathrm{A}) = \mathrm{dh}(\mathrm{A}) = +\infty\). Consequently, the functions \(\mathfrak{p} \mapsto \mathrm{di}_{A_{\mathfrak{p}}}(\mathrm{N}_{\mathfrak{p}})\) and \(\mathfrak{p} \mapsto \mathrm{dh}(\mathrm{A}_{\mathfrak{p}})\) are not in general upper semicontinuous.

